% !TEX program = xelatex
% !TEX builder = latexmk
% !Mode:: "TeX:UTF-8:Hard"
%%%%%%%%%%%%%%%%%%%%%%%%%%%%%%%%%%%%%%%%%%%%%%%%%%%%%%%%%%%%%
%% BAC01
%%
%% Andreas Landgraf (2015)
%%
%% see "How to Write a Thesis in LaTeX"
%%     https://www.sharelatex.com/blog/2013/08/02/thesis-series-pt1.html
%%%%%%%%%%%%%%%%%%%%%%%%%%%%%%%%%%%%%%%%%%%%%%%%%%%%%%%%%%%%%

\def\campusSubject{BACHELOR THESIS}
\def\campusTitle{Analysis and Comparison of Wireless Personal Area Networking Technologies for low-power Smart Home Applications}
\def\campusAuthor{Andreas Landgraf}
\def\campusPersKnz{1310317046}
\def\campusCity{Graz}
\def\campusDate{\today}
\def\campusKeywords{WPAN, Smart Home, Bluetooth, Bluetooth Smart, UWB, ZigBee, 6LoWPAN, Wi-Fi, 802.11ad}

\documentclass[
  headexclude,
  footexclude,
  10pt,
  a4paper,
  oneside,
  captions=tableheading,
  bibliography=totoc,
  listof=totoc
]{scrbook}

\RequirePackage[
  automark,
  headsepline
]{scrlayer-scrpage}

\RequirePackage{geometry}
\RequirePackage{layouts}
\RequirePackage{calc}
\RequirePackage{setspace}
\RequirePackage{titlesec}
\RequirePackage[austrian,american]{babel}
\RequirePackage[utf8]{inputenc}
\RequirePackage[T1]{fontenc}
\RequirePackage{fontspec}
\RequirePackage[hyphens]{url}
\RequirePackage{csquotes}
\RequirePackage{amsmath,amsthm,amssymb}
\RequirePackage{color}
\RequirePackage[dvipsnames,svgnames,table]{xcolor}
\RequirePackage{booktabs}
\RequirePackage{textcomp}
\RequirePackage{graphicx}
\RequirePackage{tabularx}
\RequirePackage{multirow}
\RequirePackage[font=small,format=hang,labelfont=bf]{caption}
\RequirePackage{subcaption}
\RequirePackage[colorinlistoftodos,obeyDraft]{todonotes}
\RequirePackage[toc,acronym,style=long,nonumberlist]{glossaries}

\RequirePackage[
  backend=biber,
  style=campus02,
  natbib=true,
  sorting=nty,
  sortlocale=en_US,
  autocite=footnote,
  related=true,
  backref=false,
  hyperref=auto,
  isbn=false,
  url=false,
  citereset=none,
  maxcitenames=3,
  maxbibnames=100,
  block=none,
  citecounter=true,
]{biblatex}

\RequirePackage{enumitem}

\RequirePackage{footnote}
\RequirePackage[multiple]{footmisc}
\RequirePackage{chngcntr}

\RequirePackage[
  pdftitle={\campusTitle},
  pdfauthor={\campusAuthor},
  pdfsubject={\campusSubject, \campusDate},
  pdfkeywords={\campusKeywords},
  pdffitwindow=true,
  pdfpagelayout={OneColumn},
  bookmarks=true,
  bookmarksnumbered=true,
  bookmarksopen=true,
  hyperindex=true,
  colorlinks=false,
  linkcolor=black,
  urlcolor=black,
  citecolor=black,
  unicode=true,
  hidelinks,
  hyperfootnotes=false
]{hyperref}
\RequirePackage{nameref}
\RequirePackage{cleveref}

%%%%%%%%%%%%%%%%%%%%%%%%%%%%%%%%%%%%%%%%%%%%%%%%%%%%%%%%%%%%%
% Bibliographies
%%%%%%%%%%%%%%%%%%%%%%%%%%%%%%%%%%%%%%%%%%%%%%%%%%%%%%%%%%%%%

\counterwithout{footnote}{chapter}

\DeclareLanguageMapping{american}{american-apa}
\NewBibliographyString{retrieved}
\NewBibliographyString{from}
\DefineBibliographyStrings{american}{
  retrieved = {},
  from = {},
}

\addbibresource{literature/mendeley.bib}
\addbibresource{literature/internet.bib}

\defbibfilter{printed}{
  \type{book} \or
  \type{phdthesis} \or
  \type{article} \or
  \type{incollection} \or
  \( \type{techreport} \and \not \keyword{RFC4944} \) \or
  \type{inproceedings}
}

\defbibfilter{online}{
  \type{online} \or
  \( \type{techreport} \and \keyword{RFC4944} \)
}

\newenvironment{bibhyphenation}
{\hyphenpenalty=20
  \tolerance=9999
  \exhyphenpenalty=0
  \linepenalty=0
  \setlength{\emergencystretch}{3em}}
{}

\crefname{table}{table}{tables}
\Crefname{table}{Table}{Tables}
\crefname{figure}{figure}{figures}
\Crefname{figure}{Figure}{Figures}
\crefname{section}{section}{sections}
\Crefname{section}{Section}{Sections}
\crefname{chapter}{chapter}{chapters}
\Crefname{chapter}{Chapter}{Chapters}

%%%%%%%%%%%%%%%%%%%%%%%%%%%%%%%%%%%%%%%%%%%%%%%%%%%%%%%%%%%%%

\geometry{
  left=2.5cm,
  right=2cm,
  top=2.5cm,
  bottom=2.2cm,
}

\graphicspath{{images/}}
\setmainfont{Arial}

\makeglossaries

% Acronyms
\newacronym{wban}{WBAN}{Wireless Body Area Network}
\newacronym{wpan}{WPAN}{Wireless Personal Area Network}
\newacronym{wlan}{WLAN}{Wireless Local Area Network}
\newacronym{wcan}{WCAN}{Wireless Campus Area Network}
\newacronym{wman}{WMAN}{Wireless Metropolitan Area Network}
\newacronym{wwan}{WWAN}{Wireless Wide Area Network}
\newacronym{wran}{WRAN}{Wireless Regional Area Network}
\newacronym{wgan}{WGAN}{Wireless Global Area Network}
\newacronym{han}{HAN}{Home Area Network}
\newacronym{ict}{ICT}{Information and Communications Technology}
\newacronym{iot}{IoT}{Internet of Things}
\newacronym{itu}{ITU}{International Telecommunication Union}
\newacronym{rtd}{RTD}{Research and Technical Development}
\newacronym{sms}{SMS}{Social Media Services}
\newacronym{usb}{USB}{Universal Serial Bus}
\newacronym{wusb}{WUSB}{Certified Wireless USB}
\newacronym{nb}{NB}{Narrowband}
\newacronym{uwb}{UWB}{Ultra Wideband}
\newacronym{hbc}{HBC}{Human Body Communications}
\newacronym{rfid}{RFID}{Radio Frequency Identification}
\newacronym{nfc}{NFC}{Near Field Communication}
\newacronym{pc}{PC}{Personal Computer}
\newacronym{pda}{PDA}{Personal Digital Assistant}
\newacronym{ism}{ISM}{Industrial, Scientific and Medical}
\newacronym{siso}{SISO}{Single-Input Single-Output}
\newacronym{mimo}{MIMO}{Multiple-Input Multiple-Output}
\newacronym{mumimo}{MU-MIMO}{Multi-User MIMO}
\newacronym{dsss}{DSSS}{Direct Sequence Spread Spectrum}
\newacronym{ofdm}{OFDM}{Orthogonal Frequency-Division Multiplexing}
\newacronym{wimax}{WiMAX}{Worldwide Interoperability for Microwave Access}
\newacronym{bwa}{BWA}{Broadband wireless access}
\newacronym{gsm}{GSM}{Global System for Mobile Communications}
\newacronym{tdma}{TDMA}{Time Division Multiple Access}
\newacronym{cdma}{CDMA}{Code Division Multiple Access}
\newacronym{gprs}{GPRS}{General Packet Radio Service}
\newacronym{edge}{EDGE}{Enhanced Data Rates for GSM Evolution}
\newacronym{umts}{UMTS}{Universal Mobile Telecommunications System}
\newacronym{hspa+}{HSPA+}{Evolved High-Speed Packet Access}
\newacronym{wwww}{WWWW}{World Wide Wireless Web}
\newacronym{http}{HTTP}{Hypertext Transfer Protocol}
\newacronym{when}{WHEN}{Wireless Home Entertainment Network}
\newacronym{whan}{WHAN}{Wireless Home Automation Network}
\newacronym{rf}{RF}{Radio Frequency}
\newacronym{m2m}{M2M}{Machine-to-Machine}
\newacronym{h2m}{H2M}{Human-to-Machine}
\newacronym{mih}{MiH}{Machine-in-Human}
\newacronym{whn}{WHN}{Wireless Home Network}
\newacronym{mac}{MAC}{Media Access Control}
\newacronym{l2cap}{L2CAP}{Logical Link Control and Adaption Protocol}
\newacronym{rfcomm}{RFCOMM}{Radio Frequency Communication}
\newacronym{tcs}{TCS}{Telephony Control Protocol Specification}
\newacronym{sdp}{SDP}{Service Discovery Protocol}
\newacronym{hci}{HCI}{Host Control Interface}
\newacronym{fhss}{FHSS}{Frequency Hopping Spread Spectrum}
\newacronym{gfsk}{GFSK}{Gaussian Frequency-Shift Keying}
\newacronym{fsk}{FSK}{Frequency-Shift Keying}
\newacronym{ble}{BLE}{Bluetooth Low Energy}
\newacronym{att}{ATT}{Attribute Protocol}
\newacronym{smp}{SMP}{Security Manager Protocol}
\newacronym{gatt}{GATT}{Generic Attribute Profile}
\newacronym{gap}{GAP}{Generic Access Profile}
\newacronym{pdu}{PDU}{Protocol Data Unit}
\newacronym{crc}{CRC}{Cyclic Redundancy Check}
\newacronym{single_uwb}{DS UWB}{Direct Sequence UWB}
\newacronym{multi_uwb}{MB-OFDM UWB}{Multiband Orthogonal Frequency Division Multiplexing UWB}
\newacronym{ffd}{FFD}{Full-Function Device}
\newacronym{rfd}{RFD}{Reduced-Function Device}
\newacronym{zdo}{ZDO}{ZigBee Device Object}
\newacronym{aps}{APS}{Application Support SubLayer}
\newacronym{6lowpan}{6LoWPAN}{IPv6 over Low-Power Wireless Personal Area Network}
\newacronym{lowpan}{LoWPAN}{Low-Power Wireless Personal Area Network}
\newacronym{aes}{AES}{Advanced Encryption Standard}
\newacronym{pan}{PAN}{Personal Area Network}
\newacronym{ip}{IP}{Internet Protocol}
\newacronym{ipv6}{IPv6}{Internet Protocol version 6}
\newacronym{udp}{UDP}{User Datagram Protocol}
\newacronym{icmp}{ICMP}{Internet Control Message Protocol}
\newacronym{tcp}{TCP}{Transmission Control Protocol}
\newacronym{pbss}{PBSS}{Personal Basic Service Set}
\newacronym{dsl}{DSL}{Digital Subscriber Line}
\newacronym{rfc}{RFC}{Request for Comments}

\renewcommand*{\glspostdescription}{}

\setcounter{secnumdepth}{3}
\setcounter{tocdepth}{3}

\clubpenalty = 10000
\widowpenalty = 10000
\displaywidowpenalty = 10000

\frenchspacing
\renewcommand{\singlespacing}{\setstretch{1.09}}
\renewcommand{\onehalfspacing}{\setstretch{1.45}}
\singlespacing
\newlength{\singlebaselineskip}
\setlength{\singlebaselineskip}{\baselineskip}
\onehalfspacing
\renewcommand{\arraystretch}{1.09}
\addtokomafont{sectioning}{\setstretch{1}\selectfont}

\renewcommand*{\chapterheadstartvskip}{\vspace*{4mm}}
\renewcommand*{\chapterheadendvskip}{\vspace*{0mm}}

\titlespacing*{\section}{0mm}{4mm}{0mm}
\titlespacing*{\subsection}{0mm}{4mm}{0mm}

\setlength{\parskip}{3mm}
\setlength{\itemsep}{1em}

\renewcommand{\chaptermark}[1]{\markboth{#1}{}}

% clear header and footer
\clearscrheadfoot
\ihead[]{\leftmark}
\chead[]{}
\ohead[]{}
\ifoot[]{}
\cfoot[]{}
\ofoot[]{\thepage}

\setkomafont{title}{\fontspec{Arial}\Large\bfseries}
\setkomafont{pageheadfoot}{\fontspec{Arial}\normalsize}
\setkomafont{chapter}{\fontspec{Arial}\Large\bfseries}
\setkomafont{section}{\fontspec{Arial}\Large\bfseries}
\setkomafont{subsection}{\fontspec{Arial}\large\bfseries}
\setkomafont{chapterentry}{\fontspec{Arial}\bfseries}
\addtokomafont{chapter}{\MakeUppercase}

%%%%%%%%%%%%%%%%%%%%%%%%%%%%%%%%%%%%%%%%%%%%%%%%%%%%%%%%%%%%%

\renewcommand*{\maketitle}{
  \clearpage
  \thispagestyle{empty}
  \null
  \begin{center}
    {\textbf{\campusSubject} \par}
    \vskip 2.5em
    {\hrule\vskip 2em \titlefont\huge\campusTitle\vspace*{1em}\hrule\vspace*{1.5em}\par}
    {accomplished at \par}
    \vskip .6em
    {\parbox[c]{6.6cm}{\includegraphics[width=6.6cm]{images/Campus02_Logo.pdf}} \par}
    \vskip .1em
    {FACHHOCHSCHULE DER WIRTSCHAFT \par}
    \vskip 2em
    {Bachelor Degree Programme \\ Innovation Management \par}
    \vskip 2em
    {by \\ {\Large \textbf{\campusAuthor}} \\ \campusPersKnz \par}
    \vskip 2em
    {supervised and reviewed by \par}
    {\normalsize Dipl.-Ing. Manuela Reinisch \\ FH-Prof. Dipl.-Ing. Dr.techn. Michael Terler \par}
    \vfill
    {
      \begin{flushleft}
        \hfill\parbox[c]{7.3cm}{\includegraphics[width=5cm]{images/Signature}}
        \begin{tabular}{lclcl}
          {\campusCity, \campusDate} & ~\hfill & \hspace*{3cm} & ~\hfill &  \underline{\hspace*{5cm}} \\
          ~\hfill & ~ &   & ~ &  \multicolumn{1}{c}{Signature}   \\
        \end{tabular}
      \end{flushleft}
    }
  \end{center}
  \par
}


\begin{document}

\frontmatter

\renewcommand*{\chapterpagestyle}{scrplain}
\pagestyle{scrplain}

% allow more flexible whitespaces to avoid hyphenation and overfull hboxes
\sloppy

\begin{titlepage}

\subject{\campusSubject}
\title{\campusTitle}

\author{\bfseries \campusAuthor}
\date{\campusCity, \campusDate}

\maketitle

\end{titlepage}
\newpage

\chapter*{Ehrenwörtliche Erklärung}

\begin{otherlanguage}{austrian}
Ich erkläre ehrenwörtlich, dass ich die vorliegende Arbeit selbstständig und ohne fremde Hilfe verfasst, andere als die angegebenen Quellen nicht benützt und die benutzten Quellen wörtlich zitiert sowie inhaltlich entnommene Stellen als solche kenntlich gemacht habe.
\end{otherlanguage}

{
  \begin{flushleft}
    \vskip 2em
    \hfill\parbox[c]{7.3cm}{\includegraphics[width=5cm]{images/Signature}}
    \begin{tabular}{lclcl}
      {\campusCity, \campusDate} & ~\hfill & \hspace*{3cm} & ~\hfill &  \underline{\hspace*{5cm}} \\
      ~\hfill & ~ &   & ~ &  \multicolumn{1}{c}{Signature}   \\
    \end{tabular}
  \end{flushleft}
}

%\chapter*{Acknowledgement}

%Firstly, I would like to express my sincere gratitude to the supernovae for the creation of all substantial atoms and molecules our fantastic world consists of. Their ephemeral power made it possible to form all known and unknown creatures and plants. Without this furtune it would not be possible to live this wunderful life. I thank my family for supporting me spiritually throughout my life and during writing this thesis.

\chapter*{Abstract}

Wireless networking technologies are widely used in communication devices and services in almost every area of daily business. Medical environments, security authorities, and other organizations use them to increase their performance. In particular the Smart Home sector is one of the areas that has been researched extensively. From ecological and technological point of view, the opportunities for new technologies are vast. Therefore, it is important to compare such technologies using ecological parameters. This thesis will give an overview of the wireless network types, ranging from body area to global area networks, and introduces the different Smart Home networks. Furthermore, the major wireless networking technologies of personal area networks, Bluetooth, Bluetooth Smart, Ultra Wideband, ZigBee, Internet Protocol version 6 over Low-Power Wireless Personal Area Network and Wi-Fi, are discussed in detail. A comparison is made with current chipset manufacturers implementations to create a tabular overview of these technologies and their suitability in Smart Home networks. Results show that every technology has its optimal field of application in a modern Smart Home. However, it can be concluded that further experiments will show more accurate results.

\vspace{2cm}

\textbf{Keywords:} \campusKeywords

\chapter*{Zusammenfassung}

\begin{otherlanguage}{austrian}
Drahtlose Netzwerktechnologien zählen zum Rückgrad der Kommunikationseinrichtungen unseres täglichen Arbeitsalltags. In medizinischen Einrichtungen, Sicherheitsbehördern und anderen Organisationen werden diese Technologien zur Leistungssteigerung und zur Erhöhung der Effizienz eingesetzt. Zunehmend wird in letzter Zeit auch der Smart Home Sektor Ziel intensiver Forschung im Bereich der Netzwerktechnologien. Aus ökologischer und technologischer Sicht ergeben sich Chancen für neue effizientere Technologien. Daher ist es wichtig, solche Technologien unter den Gesichtspunkten ökologischer Parameter zu vergleichen. Diese Arbeit gibt einen Überblick über die drahtlosen Netzwerktypen, die von körpernahen bis hin zu globalen Netzwerken reichen und erläutert die verschiedenen Smart Home Netzwerke. Aus dem Bereich der persönlichen Netzwerke werden die wichtigsten drahtlosen Netzwerktechnologien vorgestellt. Dabei werden Bluetooth, Bluetooth Smart, Ultra-Breitband, ZigBee, Internet Protocol Version 6 für Wireless Personal Area Network mit niedrigem Energieverbrauch und Wi-Fi ausführlich diskutiert. Ein Vergleich mit Implementierungen von Chip-Herstellern führt zu einer tabellarischen Übersicht dieser Technologien und schafft die Grundlage für die Beurteilung der Eignung in Smart Home Netzwerken. Die Ergebnisse zeigen, dass jede Technologie ihr optimales Einsatzgebiet in einem modernen Smart Home hat. Es kann jedoch der Schluss gezogen werden, dass künftige Experimente zur Eignung dieser Technologien im Smart Home diese Beurteilung vervollständigen können.
\end{otherlanguage}

\vspace{2cm}

\textbf{Keywords:} \campusKeywords

\mainmatter

\pagenumbering{roman}
\setcounter{page}{1}

\pagestyle{scrheadings}
\setlength{\parindent}{0em}
\renewcommand*{\chapterpagestyle}{scrheadings}

\tableofcontents

\clearpage

\pagenumbering{arabic}
\setcounter{page}{1}

\chapter{Introduction}
\label{ch:intro}

This \LaTeX{} template is designed for the creation of thesis documents (bachelor, master, phd) and targets both beginner and experienced users of \LaTeX{}. It supports all basic functionality and requirements of a technical document such as the inclusion of graphics, math, tables, references, bibliography and much more. In contrast to a standard LaTeX document this template not only loads all state of the art packages (\path{scientific.cls}) to provide the best functions for each task, but also includes a separate document for the style/layout of the document (\path{style.tex}). It therefore tries to separate functionallity and layout as much as possible. And the best, everything is documented in the code.

\chapter{Theory}
\label{ch:theory}

Duis porta orci. Integer eu arcu at enim tempus facilisis. Pellentesque dignissim orci sed est. Etiam elementum laoreet mi. Donec nunc sapien, dictum in, tristique sed, aliquam vitae, massa. Morbi magna magna, vestibulum tempor, lobortis non, convallis nec, nibh. In sed nibh. Suspendisse adipiscing dictum pede. Suspendisse non augue. Lorem ipsum dolor sit amet, consectetuer adipiscing elit. Pellentesque lacinia, velit sed commodo convallis, diam dolor consequat ligula, a scelerisque quam neque et purus. Praesent vel augue. Sed lectus leo, dignissim eget, vulputate eu, auctor ut, nulla. Vivamus a quam. Nulla tellus. Pellentesque tempor pulvinar nunc.

\section{Section heading}

Fusce vitae quam eu lacus pulvinar vulputate. Suspendisse potenti. Aliquam imperdiet ornare nibh. Cras molestie tortor non erat. Donec dapibus diam sed mauris laoreet volutpat. 

\subsection{Subsection heading}

Aliquam dignissim laoreet mi. Duis pulvinar nulla id velit euismod fringilla. In hac habitasse platea dictumst. Vestibulum dolor tellus, gravida a, condimentum nec, laoreet ut, nisi. Aenean lacus purus, tristique in, sagittis sit amet, pellentesque non, neque. Sed egestas nibh vitae velit. Nunc adipiscing. Donec sed lectus. Donec ultrices lacus nec orci. Fusce sit amet nulla. Suspendisse vulputate, mi nec nonummy sodales, ligula massa molestie est, at sagittis nisi est in leo.

\subsubsection{Subsubsection heading}

Mauris rutrum volutpat massa. Suspendisse potenti. Nam varius. Fusce nec leo. Morbi vestibulum augue ac justo. Vivamus in odio in turpis pharetra blandit. Mauris aliquet ullamcorper libero. Integer quam mi, venenatis ut, tristique ut, tempus at, ipsum. Donec malesuada. In quis tellus et ipsum hendrerit imperdiet. Vivamus sapien ipsum, suscipit sed, gravida a, lacinia laoreet, ipsum. Quisque augue. Nulla justo enim, auctor at, hendrerit nec, porttitor non, urna. Duis tincidunt tincidunt dui. Lorem ipsum dolor sit amet, consectetuer adipiscing elit. Suspendisse potenti. Aenean sit amet mauris. In luctus purus nec lorem. Proin orci tortor, tempus sit amet, molestie hendrerit, placerat egestas, tellus.


\chapter{Experiments}
\label{ch:experiments}

\gls{MFD}
\gls{CPA}
\gls{NA}
\gls{MMI}
\gls{SLM}
\gls{LCD}
\gls{px}
\gls{DNA}
\gls{DOF}
\gls{PSF}
\gls{SNOM}
\gls{FWHM}


\chapter{Results}
\label{ch:results}

\begin{figure}[ht]
	\centering
	\missingfigure[figwidth=6cm]{Testing a long text string}
	\caption{Typical figure, source: Own diagram}
	\label{fig:1}
\end{figure}

\chapter{Key Parameters}
\label{ch:key-params}

In ubiquitous computing of Smart Home networks, the different wireless networking protocols do not perform well in any Smart Home scenario. Therefore, the wireless networking technology which performs best in the specific Smart Home application should be choosen. To perform a comparison of wireless networking technologies in accordance to different Smart Home applications, a wireless protocol should address the following issues: \footcite[Cf.][7289--7293]{Mendes2015}

\begin{itemize}
  \item Energy efficiency: The low power consumption of the transmission layers.
  \item Bandwidth efficiency: Spread spectrum efficiency, channel width, and coexistence stability.
  \item Scalability: The number of nodes supported in a network.
  \item Transmission Range: The theoretical coverage range and the penetration of obstacles.
  \item Throughput: The maximum rate of transmission.
  \item Latency: Theoretical latency between two nodes.
  \item Mobility: Efficient mobile node management.
\end{itemize}

The taxonomy of parameters for the wireless networking technologies is shown in \cref{fig:taxonomy_of_parameters}. The evaluation of the parameters is discussed in \cref{ch:eval}.

\begin{figure}[ht]
  \centering
  \includegraphics[
    width=12cm,
  ]{images/taxonomy_of_parameters}
  \caption{Taxonomy of parameters for evaluation of wireless networking technologies, adapted from: \cite[7290]{Mendes2015}}
  \label{fig:taxonomy_of_parameters}
\end{figure}

\chapter{Evaluation and Comparison}
\label{ch:eval}

\begin{table}[h!]
  \renewcommand{\arraystretch}{2}
  \newcolumntype{P}[1]{>{\raggedright}p{#1}}
  \centering
  \begin{tabular}{ P{2.3cm} P{1.8cm} P{1.8cm} P{1.8cm} P{2cm} P{2cm} >{\arraybackslash}P{1.8cm} }
    \toprule
    \rowcolor{Gainsboro}
    \textbf{Parameters} & \textbf{Bluetooth} & \textbf{Bluetooth Smart} & \textbf{UWB} & \textbf{ZigBee} & \textbf{6LoWPAN} & \textbf{Wi-Fi} 802.11ad \\
    \midrule
    \textbf{Chipset} & CC2560 \footnotemark & CC2640 \footnotemark & {PL3120 \footnotemark \\ PL3130 \footnotemark} & CC2630 \footnotemark & CC2630 \footnotemark & Wil6200 \footnotemark \\
    \textbf{Frequency Band} & 2.4~GHz to 2.48~GHz & 2.4~GHz to 2.48~GHz & 3.1~GHz to 5.8~GHz & 868.3~MHz / 902~MHz to 928~MHz / 2.4~GHz to 2.48~GHz & 868.3~MHz / 902~MHz to 928~MHz / 2.4~GHz to 2.48~GHz & 57.24~GHz to 65.88~GHz \\
    \textbf{Number of Channels} & 79 & 40 & 3 & 27 & 27 & 4 \\
    \textbf{Channel Bandwidth} & 1~MHz & 2~MHz & 528~MHz & 1~MHz / 2~MHz / 5~MHz & 1~MHz / 2~MHz / 5~MHz & 2.16~GHz \\
    \textbf{Modulation} & \gls{gfsk} & \gls{gfsk} & \gls{ofdm} & \gls{dsss} & \gls{dsss} & \gls{ofdm} \\
    \textbf{Bandwidth Spreading} & \gls{fhss} & \gls{fhss} & \gls{dsss} & \gls{dsss} & \gls{dsss} & \gls{ofdm} \\
    \textbf{Max. Data Rate} & 2.1~Mbps & 0.27~Mbps & 675~Mbps & 20~kbpss / 40~kbps / 250~kbps & 20~kbps / 40~kbps / 250~kbps & 6.75~Gbps \\
    \textbf{Transmission Range} & 10~m & 100~m & 10~m & 10~m & 10~m & 10~m \\
    \textbf{Energy Consumption} & 1~W & 500~mW & 330~mW & 20~mW & 20~mW & 700~mW \\
    \textbf{Network Type} & Piconet / Scatternet & Scatternet & Piconet & Scatternet & Piconet & Peer-to-Peer \\
    \textbf{Max. Number of Nodes} & 7 & 5917 & 16 & 1024 & 100 & 1 \\
    \textbf{Latency} & 100~ms& 6~ms & <1~ms & 100~ms & 100~ms & 10~{\textmu s} \\
    \textbf{Mobility} & Yes & Yes & Yes & Yes & Yes & No \\
    \bottomrule
  \end{tabular}
  \caption{Parameters of \gls{wpan} technologies, source: Own diagram}
  \label{table:parameters}
  \vspace{-1.5em}
\end{table}
\footcitetext[Cf.]{cc2560}
\footcitetext[Cf.]{cc2640}
\footcitetext[Cf.]{pl3120}
\footcitetext[Cf.]{pl3130}
\footcitetext[Cf.]{cc2630}
\footcitetext[Cf.]{cc2630}
\footcitetext[Cf.]{wil6200}

The parameters of the wireless networking technologies, shown in \cref{table:parameters}, were researched from chipset manufacturers and combined with the information of \cref{ch:wpan}. Furthermore, the following comparison can be conducted as illustrated in \cref{table:comparison}.

The applicability of the Bluetooth technology is far away from the expectations. Because of the low data rate, Bluetooth is not usable for \gls{when} applications. The reason for the ordinary rating in the other \gls{whn} categories, is due to the small number of network nodes.

A totally different picture shows Bluetooth Smart. Caused to the low data rate, the Bluetooth Smart technology is not well suitable for \gls{when} applications. But Bluetooth Smart is a good choice for \gls{whan} and wireless \gls{m2m} networks. Both network types have the requirements of very low power consumption, a scalable network architecture and the support for hundreds a network nodes. Bluetooth Smart offer a good coverage range and is also applicable for outdoor sensors. But Bluetooth Smart is not very easily integratable into heterogenous \gls{whn} environments, due to the non-\gls{ip} stack.

\begin{table}[h!]
  \renewcommand{\arraystretch}{2}
  \newcolumntype{P}[1]{>{\raggedright}p{#1}}
  \centering
  \begin{tabular}{ P{3cm} P{2.5cm} P{2.5cm} P{2.5cm} >{\arraybackslash}P{2.5cm} }
    \toprule
    \rowcolor{Gainsboro}
    \textbf{Wireless Networking Technologies} & \textbf{Wireless Home Entertainment Network} & \textbf{Wireless Home Automation Network} & \textbf{Wireless \gls{m2m} Network} & \textbf{Wireless Home Network} \\
    \midrule
    \textbf{Bluetooth} & Low & Medium & Medium & Medium \\
    \textbf{Bluetooth Smart} & Low & High & High & Medium \\
    \textbf{UWB} & High & Medium & Medium & Low \\
    \textbf{ZigBee} & Low & High & High & Medium \\
    \textbf{6LoWPAN} & Low & High & High & High \\
    \textbf{Wi-Fi 802.11ad} & High & Low & Low & Medium \\
    \bottomrule
  \end{tabular}
  \caption{Aptitude of \gls{wpan} technologies in Smart Home networks, source: Own diagram}
  \label{table:comparison}
  \vspace{-1.5em}
\end{table}

The \gls{uwb} technology is very well suited for \gls{when}. \gls{uwb} offer a high data rate of several hundreds of Mbps. The technology is very robust to interference and has low power consumption. \gls{uwb} is a versatile technology, just the number of network nodes is not large enough to be the first choice in \gls{whan} and wireless \gls{m2m} networks. The non-\gls{ip} stack is also a barrier for heterogenous \gls{whn} environments.

ZigBee is the technology for the \gls{iot}, therefore the rating is high in the respective categories, \gls{whan} and wireless \gls{m2m} networks. Low power consumption and a large number of network nodes are the main requirements for those network types. ZigBee does not offer high data rates, nor is the protocol not \gls{ip} conform. Therefore, ZigBee is not a good choice for \gls{when} or heterogenous \gls{whn} environments.

The next candidate makes that possible, \gls{6lowpan}: Easy integratable into \gls{ip} networks, low power consumption and a high maximum number of network nodes. \gls{6lowpan} is well suited for \gls{whan}, wireless \gls{m2m} networks and heterogenous \gls{whn} environments. Because of the low data rate, \gls{6lowpan} is not very attractive to \gls{when}.

The escapee from the Wi-Fi technologies, the IEEE 802.11ad standard, is very well suited for high speed data transmissions as needed in \glspl{when}. But, the number of simultaneous data transfers to different network nodes is very limited. That is why the IEEE 802.11ad is not applicable in \gls{whan} or wireless \gls{m2m} networks. Due to the fact that IEEE 802.11ad \gls{mac} layer is combinable with other IEEE 802.11 standards in multi-\gls{mac} chipsets, the IEEE 802.11ad is integratable in heterogenous \gls{whn} environments.

\chapter{Conclusion and Future Work}
\label{ch:conclusion}

\begin{table}[h!]
	\centering
	\begin{tabular}{||c c c c||}
    \hline
    Col1 & Col2 & Col2 & Col3 \\ [0.5ex] 
    \hline\hline
    1 & 6 & 87837 & 787 \\ 
    2 & 7 & 78 & 5415 \\
    3 & 545 & 778 & 7507 \\
    4 & 545 & 18744 & 7560 \\
    5 & 88 & 788 & 6344 \\ [1ex]
    \hline
    \end{tabular}
	\caption{Typical table example}
	\label{table:1}
\end{table}


% uncomment to include all literature sources
%\nocite{*}

\clearpage

\ihead[]{References}
\printbibheading[heading=bibintoc, title=References]
\begin{bibhyphenation}
  \printbibliography[heading=subbibintoc, title=Bibliography, filter={printed}]
  \printbibliography[heading=subbibintoc, title=Webography, filter={online}]
\end{bibhyphenation}

\clearpage

\ihead[]{\leftmark}
\listoffigures
\listoftables
\printglossary[type=\acronymtype]

\end{document}